% GENERATED FILE. Edit canonical lessons, then run npm run generate:books.
% canonical-source-hash: fnv1a64:f96c04e2f5ea1a7a
% canonical-lessons: GU-C176-ghadiyal, GU-C176-sada, GU-C176-sava, GU-C176-ponu, GU-C176-lagbhag, GU-R176-review-171-176

\chapter{Hours and quarters}
\label{ch:gu-hours-and-quarters}
\hlchaptermodality{176}

\begin{chapteropening}
\textbf{By the end of this chapter:} \emph{I can say how long something will take in hours and quarters.}

Complete the last lesson of chapter 176: I can say how long something will take in hours and quarters.
\end{chapteropening}

\section[\texorpdfstring{\gu{ઘડિયાળ} (ghaḍiyāḷ)}{ઘડિયાળ (ghaḍiyāḷ)}]{\gu{ઘડિયાળ} (ghaḍiyāḷ) — clock, watch}

\label{lesson:GU-C176-ghadiyal}



\begin{quote}
Before the new one: say the Gujarati for half, then the Gujarati for constantly.
\end{quote}

\subsection*{You'll want to know: \gu{ઘડિયાળ}}
\textbf{\gu{ઘડિયાળ}} — \emph{ghaḍiyāḷ} — \textquotedblleft{}clock, watch\textquotedblright{}.

\begin{quote}
\textbf{\gu{હું} \gu{નવી} \gu{ઘડિયાળ} \gu{ખરીદીશ}.} — \emph{hũ navī ghaḍiyāḷ kharīdīś} — \textquotedblleft{}I will buy a new watch.\textquotedblright{}
\end{quote}

A feminine noun.

\begin{cousinweb}
\textbf{\gu{ઘડિયાળ}} (\emph{ghaḍiyāḷ}) is a tadbhava. It comes from \emph{*ghaṭītāḍa}, \textquotedblleft{}striking the time\textquotedblright{}, a reconstructed compound of \emph{ghaṭī}, \textquotedblleft{}water-clock\textquotedblright{}, and \emph{tāḍa}, \textquotedblleft{}striking\textquotedblright{}, and CDIAL glosses the Gujarati word \textquotedblleft{}clock\textquotedblright{} (CDIAL 4413).
\end{cousinweb}

\begin{grammarlens}[title={What you've built}]
One more word for everyday life at A2.
\end{grammarlens}

\subsection*{Your turn}
\begin{itemize}
\raggedright
  \item \emph{Say it:} \emph{ghaḍiyāḷ}
  \item \emph{Say it:} \emph{ghaḍiyāḷ}, once more
  \item \emph{Say it:} \emph{hũ navī ghaḍiyāḷ kharīdīś} — I will buy a new watch.
  \item \emph{Recall:} say the Gujarati for half, then the Gujarati for constantly, then say \emph{ghaḍiyāḷ} again
\end{itemize}

\begin{tcolorbox}[breakable,colback=teal!4,colframe=teal!35!black,title={Before you move on}]
What does \gu{ઘડિયાળ} mean? (\textquotedblleft{}Clock, watch\textquotedblright{}.) Say it once more.
\end{tcolorbox}
\section[\texorpdfstring{\gu{સાડા} (sāḍā)}{સાડા (sāḍā)}]{\gu{સાડા} (sāḍā) — and a half}

\label{lesson:GU-C176-sada}



\begin{quote}
Before the new one: say the Gujarati for constantly, then the Gujarati for clock.
\end{quote}

\subsection*{You'll want to know: \gu{સાડા}}
\textbf{\gu{સાડા}} — \emph{sāḍā} — \textquotedblleft{}and a half\textquotedblright{}.

\begin{quote}
\textbf{\gu{હું} \gu{સાડા} \gu{ત્રણ} \gu{કલાક} \gu{ભણીશ}.} — \emph{hũ sāḍā traṇ kalāk bhaṇīś} — \textquotedblleft{}I will study for three and a half hours.\textquotedblright{}
\end{quote}

Before a number: \textbf{\gu{સાડા} \gu{ત્રણ}}, three and a half.

\begin{cousinweb}
\textbf{\gu{સાડા}} (\emph{sāḍā}) is a tadbhava. It comes from Sanskrit \emph{sārdha}, \textquotedblleft{}plus a half\textquotedblright{}, built from \emph{sa-}, \textquotedblleft{}with\textquotedblright{}, and \emph{ardha}, \textquotedblleft{}half\textquotedblright{}; CDIAL lists the Gujarati word (CDIAL 13369).
\end{cousinweb}

\begin{grammarlens}[title={What you've built}]
One more word for everyday life at A2.
\end{grammarlens}

\subsection*{Your turn}
\begin{itemize}
\raggedright
  \item \emph{Say it:} \emph{sāḍā}
  \item \emph{Say it:} \emph{sāḍā}, once more
  \item \emph{Say it:} \emph{hũ sāḍā traṇ kalāk bhaṇīś} — I will study for three and a half hours.
  \item \emph{Recall:} say the Gujarati for constantly, then the Gujarati for clock, then say \emph{sāḍā} again
\end{itemize}

\begin{tcolorbox}[breakable,colback=teal!4,colframe=teal!35!black,title={Before you move on}]
What does \gu{સાડા} mean? (\textquotedblleft{}And a half\textquotedblright{}.) Say it once more.
\end{tcolorbox}
\section[\texorpdfstring{\gu{સવા} (savā)}{સવા (savā)}]{\gu{સવા} (savā) — and a quarter}

\label{lesson:GU-C176-sava}



\begin{quote}
Before the new one: say the Gujarati for clock, then the Gujarati for and a half.
\end{quote}

\subsection*{You'll want to know: \gu{સવા}}
\textbf{\gu{સવા}} — \emph{savā} — \textquotedblleft{}and a quarter\textquotedblright{}.

\begin{quote}
\textbf{\gu{કામ} \gu{સવા} \gu{કલાક} \gu{ચાલશે}.} — \emph{kām savā kalāk cālśe} — \textquotedblleft{}The work will take an hour and a quarter.\textquotedblright{}
\end{quote}

\begin{cousinweb}
\textbf{\gu{સવા}} (\emph{savā}) is a tadbhava. It comes from Sanskrit \emph{sapāda}, \textquotedblleft{}increased by a fourth\textquotedblright{}, built from \emph{sa-}, \textquotedblleft{}with\textquotedblright{}, and \emph{pāda}, \textquotedblleft{}foot, quarter\textquotedblright{}; CDIAL glosses the Gujarati word \textquotedblleft{}plus a quarter\textquotedblright{} (CDIAL 13134).
\end{cousinweb}

\begin{grammarlens}[title={What you've built}]
One more word for everyday life at A2.
\end{grammarlens}

\subsection*{Your turn}
\begin{itemize}
\raggedright
  \item \emph{Say it:} \emph{savā}
  \item \emph{Say it:} \emph{savā}, once more
  \item \emph{Say it:} \emph{kām savā kalāk cālśe} — The work will take an hour and a quarter.
  \item \emph{Recall:} say the Gujarati for clock, then the Gujarati for and a half, then say \emph{savā} again
\end{itemize}

\begin{tcolorbox}[breakable,colback=teal!4,colframe=teal!35!black,title={Before you move on}]
What does \gu{સવા} mean? (\textquotedblleft{}And a quarter\textquotedblright{}.) Say it once more.
\end{tcolorbox}
\section[\texorpdfstring{\gu{પોણું} (poṇũ)}{પોણું (poṇũ)}]{\gu{પોણું} (poṇũ) — three quarters}

\label{lesson:GU-C176-ponu}



\begin{quote}
Before the new one: say the Gujarati for and a half, then the Gujarati for and a quarter.
\end{quote}

\subsection*{You'll want to know: \gu{પોણું}}
\textbf{\gu{પોણું}} — \emph{poṇũ} — \textquotedblleft{}three quarters\textquotedblright{}.

\begin{quote}
\textbf{\gu{હું} \gu{પોણો} \gu{કલાક} \gu{ચાલીશ}.} — \emph{hũ poṇo kalāk cālīś} — \textquotedblleft{}I will walk for three quarters of an hour.\textquotedblright{}
\end{quote}

Before a number it means a quarter less: \textbf{\gu{પોણા} \gu{બે}}, one and three quarters.

\begin{cousinweb}
\textbf{\gu{પોણું}} (\emph{poṇũ}) is a tadbhava. It comes from \emph{*padūna}, \textquotedblleft{}less by a quarter\textquotedblright{}, a reconstructed form beside Sanskrit \emph{pādona}; CDIAL glosses the Gujarati word \textquotedblleft{}three quarters\textquotedblright{} (CDIAL 7765).
\end{cousinweb}

\begin{grammarlens}[title={What you've built}]
One more word for everyday life at A2.
\end{grammarlens}

\subsection*{Your turn}
\begin{itemize}
\raggedright
  \item \emph{Say it:} \emph{poṇũ}
  \item \emph{Say it:} \emph{poṇũ}, once more
  \item \emph{Say it:} \emph{hũ poṇo kalāk cālīś} — I will walk for three quarters of an hour.
  \item \emph{Recall:} say the Gujarati for and a half, then the Gujarati for and a quarter, then say \emph{poṇũ} again
\end{itemize}

\begin{tcolorbox}[breakable,colback=teal!4,colframe=teal!35!black,title={Before you move on}]
What does \gu{પોણું} mean? (\textquotedblleft{}Three quarters\textquotedblright{}.) Say it once more.
\end{tcolorbox}
\section[\texorpdfstring{\gu{લગભગ} (lagbhag)}{લગભગ (lagbhag)}]{\gu{લગભગ} (lagbhag) — almost, about}

\label{lesson:GU-C176-lagbhag}



\begin{quote}
Before the new one: say the Gujarati for and a quarter, then the Gujarati for three quarters.
\end{quote}

\subsection*{You'll want to know: \gu{લગભગ}}
\textbf{\gu{લગભગ}} — \emph{lagbhag} — \textquotedblleft{}almost, about\textquotedblright{}.

\begin{quote}
\textbf{\gu{તે} \gu{લગભગ} \gu{બે} \gu{કલાકમાં} \gu{આવશે}.} — \emph{te lagbhag be kalākmā̃ āvśe} — \textquotedblleft{}He will come in about two hours.\textquotedblright{}
\end{quote}

\begin{cousinweb}
The origin of \textbf{\gu{લગભગ}} (\emph{lagbhag}) is not settled. CDIAL lists no Gujarati form of it, and EtymDB has no entry for it.
\end{cousinweb}

\begin{grammarlens}[title={What you've built}]
One more word for everyday life at A2.
\end{grammarlens}

\subsection*{Your turn}
\begin{itemize}
\raggedright
  \item \emph{Say it:} \emph{lagbhag}
  \item \emph{Say it:} \emph{lagbhag}, once more
  \item \emph{Say it:} \emph{te lagbhag be kalākmā̃ āvśe} — He will come in about two hours.
  \item \emph{Recall:} say the Gujarati for and a quarter, then the Gujarati for three quarters, then say \emph{lagbhag} again
\end{itemize}

\begin{tcolorbox}[breakable,colback=teal!4,colframe=teal!35!black,title={Before you move on}]
What does \gu{લગભગ} mean? (\textquotedblleft{}Almost, about\textquotedblright{}.) Say it once more.
\end{tcolorbox}
\section[(dialogue)]{Words from chapters 171–176 — a review}

\label{lesson:GU-R176-review-171-176}



\begin{quote}
Say the Gujarati for taste and to taste.
\end{quote}

\subsection*{Your turn}
Say these aloud:

\begin{itemize}
\raggedright
  \item taste, to taste, bitter, salty, sweet
  \item bland, stale, to get burnt, orange, juice
  \item week, fortnight, next, last, whole
  \item October, December, Maha, wall calendar, beginning
  \item full moon day, new moon day, end, half, constantly
  \item clock, and a half, and a quarter, three quarters, almost
\end{itemize}

\begin{tcolorbox}[breakable,colback=teal!4,colframe=teal!35!black,title={Before you move on}]
Taste? (\textbf{\gu{સ્વાદ}}.) To taste? (\textbf{\gu{ચાખવું}}.) Bitter? (\textbf{\gu{કડવું}}.) Salty? (\textbf{\gu{ખારું}}.) Sweet? (\textbf{\gu{ગળ્યું}}.) Bland? (\textbf{\gu{ફીકું}}.) Stale? (\textbf{\gu{વાસી}}.) To get burnt? (\textbf{\gu{બળવું}}.) Orange? (\textbf{\gu{નારંગી}}.) Juice? (\textbf{\gu{રસ}}.) Week? (\textbf{\gu{અઠવાડિયું}}.) Fortnight? (\textbf{\gu{પખવાડિયું}}.) Next? (\textbf{\gu{આવતું}}.) Last? (\textbf{\gu{છેલ્લું}}.) Whole? (\textbf{\gu{આખું}}.) October? (\textbf{\gu{ઓક્ટોબર}}.) December? (\textbf{\gu{ડિસેમ્બર}}.) Maha? (\textbf{\gu{મહા}}.) Wall calendar? (\textbf{\gu{તારીખિયું}}.) Beginning? (\textbf{\gu{શરૂઆત}}.) Full moon day? (\textbf{\gu{પૂનમ}}.) New moon day? (\textbf{\gu{અમાસ}}.) End? (\textbf{\gu{અંત}}.) Half? (\textbf{\gu{અડધું}}.) Constantly? (\textbf{\gu{સતત}}.) Clock? (\textbf{\gu{ઘડિયાળ}}.) And a half? (\textbf{\gu{સાડા}}.) And a quarter? (\textbf{\gu{સવા}}.) Three quarters? (\textbf{\gu{પોણું}}.) Almost? (\textbf{\gu{લગભગ}}.)
\end{tcolorbox}
